% Kerangka berdasarkan samplepaper.tex (LLNCS v2.21). Preamble template dipertahankan;
% hanya paket untuk bahasa Indonesia, matematika, tabel, dan gambar yang ditambahkan.
% Membutuhkan llncs.cls dan splncs04.bst: buat project dari template resmi "Springer LNCS"
% di Overleaf, lalu unggah berkas-berkas ini di atasnya.
\documentclass[runningheads]{llncs}
\usepackage[T1]{fontenc}
\usepackage{graphicx}

% Package Tambahan (dari template)
\usepackage{float}
\usepackage{array}
\usepackage{algorithm}
\usepackage{algpseudocode}
\usepackage{times}
\usepackage{hyperref}
\usepackage{caption}
% Package Tambahan (untuk naskah ini); llncs tidak memuat paket AMS sendiri
\usepackage[indonesian]{babel}
\usepackage{amsmath}
\usepackage{amssymb}
\usepackage{booktabs}
\usepackage{xcolor}
\usepackage{tikz}
\usetikzlibrary{positioning,arrows.meta,fit,backgrounds}
\usepackage{pgfplots}
\pgfplotsset{compat=1.17}

% Nama bagian dalam bahasa Indonesia (dijalankan setelah class dan babel dimuat)
\AtBeginDocument{%
  \def\abstractname{Abstrak.}%
  \def\refname{Daftar Pustaka}%
  \def\figurename{Gambar}%
  \def\tablename{Tabel}%
  \def\keywordname{{\bfseries Kata kunci:}}%
}

% Penanda hasil yang BELUM dihasilkan. Jangan diganti dengan angka perkiraan.
\newcommand{\pending}{\textcolor{red}{\textbf{[belum ada]}}}
% Penanda teks atau informasi yang masih dibutuhkan dari penulis.
\newcommand{\todo}[1]{\textcolor{blue}{[\textbf{TODO:} #1]}}

\begin{document}

\title{Self-Supervised Learning dan Transfer Learning untuk Lokalisasi Root-Cause Service pada Microservice: Studi Perbandingan TS2Vec, TimeMAE, dan RUN pada RCAEval}
\titlerunning{SSL dan Transfer Learning untuk Lokalisasi Root-Cause Service}

\author{Fauzan Pradana \and Edy Rahman \and Zalzabila Rani}
\authorrunning{F. Pradana et al.}
\institute{Fakultas Ilmu Komputer, Universitas Brawijaya, Malang, Indonesia}
\maketitle

\begin{abstract}
Kompleksitas arsitektur microservice menyebabkan proses identifikasi sumber penyebab gangguan (Root Cause Localization) menjadi tantangan, terutama karena setiap service memiliki karakteristik metrik dan pola interaksi yang saling berkaitan. Pendekatan pembelajaran mesin berbasis data berlabel dapat digunakan untuk membantu proses lokalisasi, tetapi ketersediaan label gangguan sering kali terbatas. Salah satu pendekatan yang berpotensi mengatasi keterbatasan tersebut adalah Self-Supervised Learning (SSL), yang memungkinkan model mempelajari representasi data tanpa memerlukan label pada tahap prapelatihan. Penelitian ini bertujuan mengevaluasi penerapan SSL dan Transfer Learning untuk lokalisasi service penyebab gangguan pada sistem microservice, dengan membandingkan tiga pendekatan, yaitu TS2Vec, TimeMAE, dan RUN. Evaluasi dilakukan menggunakan dataset RCAEval, yang mencakup tiga sistem microservice, yaitu Online Boutique, Sock Shop, dan Train Ticket. Penelitian ini dirancang melalui dua skenario eksperimen, yaitu evaluasi dalam dataset yang sama (in-domain) dan transfer representasi antar-dataset (cross-dataset transfer). Pada tahap prapelatihan, model mempelajari pola data deret waktu tanpa menggunakan label, kemudian representasi yang diperoleh digunakan pada tahap penyesuaian model (fine-tuning) dengan variasi strategi FineLast dan FineAll. Kinerja lokalisasi dievaluasi menggunakan metrik Accuracy at k (AC@k) dan Average at 5 (Avg@5), dengan membandingkan hasil prapelatihan terhadap inisialisasi acak serta pendekatan heuristik. Penelitian ini diharapkan dapat memberikan pemahaman mengenai kemampuan representasi hasil pembelajaran mandiri dalam mendukung lokalisasi akar penyebab gangguan serta efektivitas transfer representasi antar-sistem microservice.


\keywords{Root Cause Localization \and Microservice \and Self-Supervised Learning \and Transfer Learning \and Representasi deret waktu}
\end{abstract}

\section{Pendahuluan}
Arsitektur microservice memecah sebuah aplikasi menjadi layanan (service) yang di-deploy secara mandiri dan saling memanggil. Pola ini meningkatkan skalabilitas dan kemudahan pemeliharaan, tetapi anomali pada satu service dapat merambat ke service yang bergantung padanya dan akhirnya menyebabkan kegagalan sistem secara keseluruhan \cite{lin2024run}. Ketika insiden terjadi, insinyur on-call harus menemukan service yang menjadi sumber gangguan. Tugas ini disebut Root Cause Localization (RCL), yaitu bagian dari Root Cause Analysis (RCA) \cite{sun2024art}. Penelusuran manual berjalan lambat, dan banyaknya metrik, log, serta trace yang dihasilkan sistem yang sedang berjalan membuat dukungan otomatis diperlukan \cite{pham2024rcaeval}.

Metode RCL terawasi bergantung pada insiden historis berlabel. Label tersebut langka dan mahal karena akar masalah harus dikonfirmasi oleh operator berpengalaman \cite{sun2024art}. Self-Supervised Learning (SSL) menawarkan cara memanfaatkan data pemantauan tanpa label: encoder dilatih lebih dulu dengan tugas pretext, kemudian diadaptasi dengan sedikit data berlabel. Metode SSL deret waktu terkini, seperti TS2Vec yang berbasis kontrastif \cite{yue2022ts2vec} dan TimeMAE yang berbasis masked autoencoder \cite{cheng2026timemae}, menargetkan representasi yang dapat ditransfer. TimeMAE melaporkan keunggulan pada kondisi label terbatas dan ketika encoder yang di-pre-train pada satu dataset di-fine-tune pada dataset lain, sedangkan TS2Vec menerapkan encoder yang dilatih pada satu dataset untuk deteksi anomali pada dua dataset lain. Namun, keduanya dievaluasi pada tolok ukur deret waktu umum, yaitu sebuah sampel berupa satu deret dengan label kelas, target peramalan, atau penanda anomali. RCL berbeda: satu kasus memuat banyak service yang saling bergantung, dan keluarannya adalah pilihan di antara service tersebut.

Metode analisis akar masalah RUN \cite{lin2024run} dan ART \cite{sun2024art} memodelkan lingkungan microservice secara langsung. RUN memakai backbone peramalan yang di-pre-train secara kontrastif di dalam pipeline causal discovery per kasus, dan dievaluasi pada data sintetis serta satu test bed. ART mempelajari perilaku sistem dari data normal dan dievaluasi pada dua dataset, salah satunya tidak publik. Tidak satu pun dari kedua studi itu menguji apakah apa yang dipelajari pada satu dataset RCA bermanfaat pada dataset lain.

\paragraph{Celah penelitian.} Sepengetahuan kami, apakah representasi self-supervised yang di-pre-train pada satu dataset RCA microservice dapat ditransfer ke dataset lain belum pernah diteliti. Bukti yang ada terpisah: encoder SSL divalidasi pada tolok ukur deret waktu umum, sedangkan metode RCA divalidasi satu dataset demi satu dataset. Studi transfer pada tolok ukur bersama juga membutuhkan kontrol yang sering tidak ada, yaitu encoder yang diinisialisasi acak dengan arsitektur sama untuk memisahkan manfaat pre-training dari manfaat fine-tuning terawasi, serta baseline heuristik sederhana.

\paragraph{Tujuan.} Penelitian ini mengkaji transfer representasi deret waktu self-supervised untuk lokalisasi root-cause service pada tolok ukur RCAEval \cite{pham2024rcaeval}. Tiga pendekatan representasi dibandingkan: TS2Vec (kontrastif), TimeMAE (masked autoencoding), dan adaptasi tahap pre-training RUN. RUN dalam bentuk aslinya dievaluasi sebagai baseline RCA tanpa transfer. Karena mekanisme pelatihan ketiga metode berbeda, analisis memisahkan tiga pertanyaan, dan satu protokol tunggal tidak dipaksakan kepada semua metode.

\paragraph{Pertanyaan penelitian.}
\begin{description}
\item[RQ1] Dapatkah encoder yang di-pre-train dengan SSL pada satu dataset RCAEval ditransfer ke dataset lain untuk lokalisasi root-cause service, dibandingkan dengan encoder yang diinisialisasi acak?
\item[RQ2] Bagaimana perbandingan TS2Vec, TimeMAE, dan RUN dari sisi (a) akurasi lokalisasi dan (b) manfaat yang diperoleh dari pre-training dan transfer?
\item[RQ3] Bagaimana proporsi data latih berlabel (10\%, 25\%, 50\%, 100\%) dan strategi fine-tuning (FineLast dibandingkan FineAll) memengaruhi akurasi serta manfaat pre-training?
\end{description}

\paragraph{Kontribusi.} (i) Protokol pembelajaran transfer untuk RCL tingkat service pada RCAEval dengan skenario in-domain dan lintas dataset, kontrol encoder acak dan heuristik, serta uji bootstrap berpasangan. (ii) Eksperimen kontrol yang mengganti data pasca-injeksi dengan data pra-injeksi untuk memeriksa bahwa model membaca anomali, bukan menghafal identitas service. (iii) Perbandingan tiga keluarga SSL dengan aturan yang memisahkan perbandingan metode, efek pre-training, dan efek label serta fine-tuning. (iv) Uraian tentang kondisi ketika pre-training dan transfer tidak membantu. \todo{Konfirmasi kontribusi setelah eksperimen TimeMAE dan RUN selesai; temuan utama ditambahkan ke paragraf ini hanya setelah tersedia.}

Sisa artikel disusun sebagai berikut. Bagian~\ref{sec:related} meninjau penelitian terkait, Bagian~\ref{sec:method} memaparkan metodologi, Bagian~\ref{sec:setup} dataset dan pengaturan eksperimen, Bagian~\ref{sec:results} dan~\ref{sec:discussion} hasil dan pembahasan, dan Bagian~\ref{sec:conclusion} kesimpulan.

\section{Tinjauan Pustaka}\label{sec:related}
\subsection{Root Cause Localization pada Sistem Microservice}
RCAEval adalah tolok ukur terbuka berisi 735 kasus kegagalan dari tiga sistem microservice, dilengkapi sejumlah baseline yang dapat direproduksi \cite{pham2024rcaeval}. Pendekatan berbasis struktur kausal yang lebih awal membangun graf dengan algoritma PC lalu merangking kandidat pada graf tersebut. RUN berargumen bahwa pendekatan itu mengabaikan urutan waktu dan menggantinya dengan neural Granger causal discovery \cite{lin2024run}. RUN terdiri atas tiga tahap: pre-training kontrastif backbone DLinear dengan dua potongan acak yang saling tumpang tindih dari deret yang sama sebagai pasangan positif tanpa pasangan negatif; jaringan peramalan per kasus dengan matriks atensi kausal yang dapat dipelajari, lalu diberi ambang dan dipangkas menjadi graf; serta PageRank dengan vektor personalisasi untuk merangking akar masalah dengan KPI pemicu (trigger) tertentu. RUN dievaluasi dengan HR@$k$ dan MRR pada data sintetis dan test bed Sock Shop (13 microservice).

ART adalah framework tak terawasi terpadu untuk deteksi anomali, triase kegagalan, dan RCL \cite{sun2024art}. ART menserialisasi metrik, log, dan trace menjadi jendela, memodelkan dependensi channel, temporal, dan call dengan Transformer encoder, GRU, dan GraphSAGE, dan melatih modelnya untuk memprediksi langkah waktu berikutnya pada data normal dengan loss galat kuadrat rata-rata. Saat diagnosis, selisih antara prediksi dan observasi membentuk representasi kegagalan tingkat instance dan tingkat sistem. Untuk RCL, instance dirangking berdasarkan kemiripan kosinus antara deviasi instance dan deviasi sistem, yang bertumpu pada pengamatan empiris bahwa instance akar masalah menyimpang lebih mirip dengan sistem secara keseluruhan. ART dievaluasi pada dua dataset, yaitu D1 (sistem e-commerce simulasi dengan 46 instance dan 210 kegagalan) dan D2 (sistem bank dengan 18 instance dan 133 kegagalan, tidak publik). Pembagian data dilakukan berdasarkan waktu (40\% kegagalan terakhir untuk pengujian). Metrik yang dilaporkan adalah Top1, Top3, dan AVG@5 dengan
$\mathrm{Top}K=\frac{1}{N}\sum_{i=1}^{N}\mathbb{1}[g_i\in P_{i,[1:K]}]$ dan $\mathrm{AVG@5}=\frac{1}{5}\sum_{K=1}^{5}\mathrm{Top}K$, dirata-ratakan dari lima kali pengulangan. Penulisnya mencatat dua keterbatasan yang relevan di sini: model terikat pada himpunan instance sehingga harus dilatih ulang ketika service berubah, dan kedua dataset berukuran kecil.

\subsection{Representasi Deret Waktu Self-Supervised}
TS2Vec mempelajari representasi pada tingkat semantik sembarang. Dua tampilan deret yang dipotong dan dimasking secara acak dikodekan oleh proyeksi masukan yang diikuti CNN dilated, dan pasangan positifnya adalah representasi timestamp yang sama pada kedua tampilan (contextual consistency). Loss menggabungkan kontras temporal dan kontras antar-instance, lalu diterapkan pada beberapa skala yang diperoleh melalui max pooling \cite{yue2022ts2vec}. TS2Vec dievaluasi pada 125 dataset UCR dan 29 dataset UEA untuk klasifikasi, dengan SVM dilatih di atas representasi yang dibekukan. TS2Vec juga dievaluasi pada peramalan dan deteksi anomali, dengan satu pengaturan memakai encoder yang dilatih pada FordA dan diuji pada Yahoo dan KPI. TimeMAE memotong deret menjadi sub-series yang tidak tumpang tindih dan memasking sebagian besarnya (60\% secara default). Transformer hanya mengodekan sub-series yang terlihat, sedangkan decoupled encoder berbasis cross-attention yang terpisah merepresentasikan posisi yang di-mask, sehingga tidak ada token mask yang masuk ke encoder utama. Tugas pretext-nya adalah masked codeword classification dengan tokenizer dan codebook yang dipelajari, serta masked representation regression terhadap target encoder yang diperbarui dengan momentum \cite{cheng2026timemae}. TimeMAE dievaluasi pada lima dataset klasifikasi dengan dua protokol yang disebutnya FineLast (encoder dibekukan) dan FineAll (fine-tuning penuh), dengan variasi proporsi label latih, serta pengaturan one-to-many ketika encoder yang di-pre-train pada HAR di-fine-tune pada empat dataset lainnya. Tidak satu pun dari kedua paper ini mengevaluasi lokalisasi di antara himpunan kandidat.

\subsection{Posisi dan Celah Penelitian}
Tabel~\ref{tab:related} membandingkan keempat metode dengan penelitian ini. Tiga perbedaan penting bagi penelitian ini.
\begin{enumerate}
\item \emph{Bentuk tugas.} TS2Vec dan TimeMAE menghasilkan satu representasi untuk satu deret, sedangkan RCL mengharuskan merepresentasikan setiap service dalam satu kasus lalu memilih satu di antaranya. Karena itu kami menambahkan scoring head di atas representasi per service (Bagian~\ref{sec:method}), yaitu pilihan rancangan yang tidak diambil dari kedua paper tersebut.
\item \emph{Peran label.} RUN dan ART tidak membutuhkan label, tetapi keduanya terikat pada dataset (RUN per kasus, ART per sistem) dan tidak menawarkan komponen yang dapat ditransfer dan sudah diuji. TS2Vec dan TimeMAE menghasilkan encoder yang dapat dipakai ulang, tetapi membutuhkan label untuk keputusan hilir. Studi transfer karena itu harus membandingkan kedua keluarga dengan aturan yang tidak merugikan salah satunya.
\item \emph{Bukti transfer.} TimeMAE menguji transfer antar-dataset klasifikasi dan TS2Vec antar-dataset deteksi anomali. ART dan RUN tidak menguji transfer, dan TS2Vec maupun TimeMAE belum pernah dievaluasi pada data RCA microservice. Kami mengadopsi istilah FineLast/FineAll serta rancangan fraksi label dari TimeMAE, dan menambahkan encoder yang diinisialisasi acak sebagai kontrol, analog dengan baseline berinisialisasi acak pada paper tersebut.
\end{enumerate}
Heuristik deviasi sederhana yang dipakai sebagai baseline di sini berkaitan dengan pengamatan ART bahwa deviasi anomali membawa informasi tentang akar masalah. Berbeda dari ART, heuristik kami memakai skor tetap dari robust $z$-score dan tidak dipelajari. ART tidak diikutsertakan dalam eksperimen. ART berfungsi sebagai acuan untuk definisi tugas, metrik (Top-$K$, AVG@5), dan pembahasan tentang pelatihan ulang ketika service berubah.
\begin{table}
\caption{Perbandingan metode terkait dengan penelitian ini (informasi diambil dari paper yang dikutip).}\label{tab:related}
\centering
\scriptsize
\setlength{\tabcolsep}{2.5pt}
\begin{tabular}{p{1.5cm}p{2.0cm}p{2.0cm}p{2.9cm}p{2.0cm}}
\toprule
 & Sinyal pembelajaran & Label dibutuhkan & Data yang dievaluasi & Transfer lintas dataset\\
\midrule
TS2Vec \cite{yue2022ts2vec} & Kontrastif hierarkis & Hanya untuk tugas hilir & UCR, UEA; peramalan; deteksi anomali & Terbatas (encoder dilatih pada FordA, diuji pada deteksi anomali Yahoo dan KPI), bukan RCA\\
TimeMAE \cite{cheng2026timemae} & Masked codeword dan regresi & Hanya untuk tugas hilir & 5 dataset klasifikasi & Ya (HAR $\rightarrow$ empat dataset), bukan RCA\\
RUN \cite{lin2024run} & Backbone kontrastif, graf Granger per kasus & Tidak ada & Sintetis; Sock Shop & Tidak diteliti\\
ART \cite{sun2024art} & Prediksi langkah berikutnya pada data normal & Tidak ada untuk RCL & D1 (publik), D2 (privat) & Tidak diteliti\\
Penelitian ini & TS2Vec, TimeMAE, RUN-encoder dengan scoring head & Label untuk fine-tuning & RCAEval RE1--RE3 & Ya (RE1, RE2, RE3)\\
\bottomrule
\end{tabular}
\end{table}


\section{Metodologi Penelitian}\label{sec:method}
Gambar~\ref{fig:pipeline} merangkum alur kerja. Kasus-kasus RCAEval diproses menjadi tensor berukuran tetap; encoder di-pre-train tanpa label pada dataset sumber; scoring head kemudian dilatih dengan label pada dataset target dalam salah satu dari dua mode fine-tuning; dan setiap model dievaluasi terhadap dua kontrol.
\begin{figure}
\centering
\resizebox{\textwidth}{!}{%
\begin{tikzpicture}[node distance=5mm, box/.style={draw, rounded corners, align=center, font=\scriptsize, minimum height=10mm, text width=2.1cm}, >=Stealth]
\node[box] (data) {RCAEval\\RE1, RE2, RE3};
\node[box, right=of data] (prep) {Preprocessing\\robust $z$-score,\\jendela 540\,detik};
\node[box, right=of prep] (pre) {Pre-training SSL\\(dataset sumber,\\tanpa label)};
\node[box, right=of pre] (ft) {Fine-tuning\\FineLast / FineAll\\(target, berlabel)};
\node[box, right=of ft] (ev) {Evaluasi\\AC@$k$, bootstrap};
\draw[->] (data)--(prep); \draw[->] (prep)--(pre); \draw[->] (pre)--(ft); \draw[->] (ft)--(ev);
\end{tikzpicture}}
\caption{Alur penelitian. Dua kontrol dipakai di seluruh eksperimen: encoder yang diinisialisasi acak (tanpa pre-training) dan baseline heuristik yang tidak memerlukan pelatihan.}\label{fig:pipeline}
\end{figure}


\subsection{Formulasi Masalah dan Notasi}
Sebuah \emph{kasus} $i$ adalah satu kegagalan yang diinjeksikan. Kasus ini memiliki himpunan $N_i$ service aplikasi kandidat yang diamati selama $T=540$ detik (300\,detik sebelum dan 240\,detik sesudah injeksi fault) melalui $C=5$ channel, yaitu CPU, memori, latency-50, latency-90, dan workload. Masukannya adalah tensor $\mathbf{X}_i\in\mathbb{R}^{N_i\times T\times C}$ dengan mask ketersediaan $\mathbf{M}_i\in\{0,1\}^{N_i\times C}$, dan labelnya adalah indeks $c_i\in\{1,\dots,N_i\}$ dari service akar masalah. Metode lokalisasi mengembalikan skor $s_{i,j}$ untuk setiap service kandidat $j$ dan merangking service berdasarkan skor yang menurun. Waktu injeksi diasumsikan diketahui, seperti pada evaluasi RCL di ART \cite{sun2024art}. Untuk semua metode, himpunan kandidat suatu kasus sama, dan keluarannya adalah rangking service, bukan rangking KPI. Bagian~\ref{sec:runmap} menjelaskan bagaimana hal ini dijamin untuk RUN.

\subsection{Preprocessing}
\begin{enumerate}
\item \textbf{Penyelarasan fitur.} Service aplikasi dipilih berdasarkan ketersediaan CPU. Jika suatu channel tidak tersedia untuk sebuah service, nilainya diisi 0 dan ditandai tidak tersedia pada $\mathbf{M}_i$.
\item \textbf{Pembersihan dan normalisasi.} Deret diresampling ke grid 1\,detik dengan forward dan backward fill. Untuk setiap channel, robust $z$-score dihitung dari median dan median absolute deviation (MAD) pada seluruh periode pra-injeksi. Simpangan baku dipakai sebagai cadangan ketika MAD degenerat, dan nilai dibatasi pada $[-10,10]$.
\item \textbf{Jendela dan tensor.} 300\,detik sebelum dan 240\,detik sesudah injeksi digabung menjadi $\mathbf{X}_i$.
\item \textbf{Pembagian data.} Kasus dibagi menjadi data latih, validasi, dan uji dengan rasio sekitar 60/20/20, terstratifikasi menurut sistem dan service akar masalah. Pre-training hanya memakai kasus latih; kasus validasi dipakai untuk pemilihan model dan early stopping.
\end{enumerate}

\subsection{Encoder dan Objective Pre-training}
Setiap encoder $g$ memetakan deret satu service menjadi vektor representasi dan diterapkan dengan bobot bersama pada setiap service dalam satu kasus. Objective di bawah ini adalah objective pada paper aslinya; $r_{i,t}$ menyatakan representasi timestamp $t$ dari deret $i$, tanda petik tunggal menandai tampilan kedua, dan $\mathcal{B}$ adalah satu batch pelatihan.

\paragraph{TS2Vec.} Encoder terdiri atas proyeksi masukan, modul timestamp masking (mask Bernoulli dengan $p=0{,}5$ pada vektor laten), dan CNN dilated dengan sepuluh blok residual \cite{yue2022ts2vec}. Dua tampilan yang dipotong secara acak dan saling tumpang tindih dikodekan, dan representasi timestamp pada bagian yang tumpang tindih membentuk pasangan positif. Dengan $\Omega$ sebagai himpunan timestamp yang tumpang tindih,
\begin{align}
\ell_{\mathrm{temp}}^{(i,t)} &= -\log\frac{\exp(r_{i,t}\cdot r'_{i,t})}{\sum_{t'\in\Omega}\bigl(\exp(r_{i,t}\cdot r'_{i,t'})+\mathbb{1}_{[t\neq t']}\exp(r_{i,t}\cdot r_{i,t'})\bigr)},\\
\ell_{\mathrm{inst}}^{(i,t)} &= -\log\frac{\exp(r_{i,t}\cdot r'_{i,t})}{\sum_{j\in\mathcal{B}}\bigl(\exp(r_{i,t}\cdot r'_{j,t})+\mathbb{1}_{[i\neq j]}\exp(r_{i,t}\cdot r_{j,t})\bigr)},
\end{align}
dan dual loss $\mathcal{L}_{\mathrm{dual}}$ adalah rata-rata $\ell_{\mathrm{temp}}^{(i,t)}+\ell_{\mathrm{inst}}^{(i,t)}$ pada seluruh deret dan timestamp. Hierarchical loss mengulangi perhitungan ini setelah max pooling berturut-turut (ukuran kernel 2) sepanjang waktu, lalu merata-ratakannya pada seluruh skala.

\paragraph{TimeMAE.} Deret dipotong menjadi sub-series yang tidak tumpang tindih, dan 60\% di antaranya di-mask secara acak pada setiap epoch. Transformer encoder $\mathcal{H}$ memproses sub-series yang terlihat dengan relative positional embedding. Decoupled encoder $\mathcal{F}$, yang diperoleh dengan mengganti self-attention menjadi cross-attention, menghitung representasi posisi yang di-mask dari posisi yang terlihat, sehingga tidak ada token mask yang masuk ke $\mathcal{H}$ \cite{cheng2026timemae}. Loss pre-training menggabungkan masked codeword classification $\mathcal{L}_c$, dengan target berupa codeword yang ditetapkan tokenizer dan codebook terpelajar pada tiap sub-series yang di-mask, dan masked representation regression $\mathcal{L}_r$, yaitu jarak $\ell_2$ kuadrat antara prediksi online encoder dan representasi dari target encoder yang diperbarui dengan momentum:
\begin{equation}
\mathcal{L}=\alpha\,\mathcal{L}_c+\beta\,\mathcal{L}_r .
\end{equation}
Target encoder diperbarui sebagai exponential moving average dari online encoder dengan koefisien 0,99.

\paragraph{RUN-original.} Pipeline RUN \cite{lin2024run} dijalankan tanpa modifikasi pada setiap kasus uji. Tiga tahapnya adalah (i) pre-training kontrastif backbone DLinear $f$ dengan projector $g$ dan loss tanpa pasangan negatif
\begin{equation}
\mathcal{L}_{\mathrm{con}}=-\tfrac{1}{2}\bigl(\cos(z_1,\mathrm{sg}(p_2))+\cos(\mathrm{sg}(p_1),z_2)\bigr),
\end{equation}
dengan $z=g(f(y))$, $p$ keluaran backbone yang bersesuaian, dan $\mathrm{sg}$ operator stop-gradient; (ii) jaringan peramalan per kasus dengan matriks atensi kausal terpelajar, diberi ambang 0,5 dan dibersihkan dari siklus dengan menghapus sisi yang korelasi Pearson-nya terendah; dan (iii) PageRank dengan vektor personalisasi, dengan skor seri diselesaikan melalui jarak akses dari titik pemicu. RUN-original tidak membutuhkan label dan tidak memiliki parameter yang dapat ditransfer, sehingga ia menjadi baseline RCA tanpa pembelajaran transfer.

\paragraph{RUN-encoder (varian adaptasi, bukan RUN asli).} RUN-encoder hanya mempertahankan tahap (i): backbone DLinear yang di-pre-train dengan $\mathcal{L}_{\mathrm{con}}$ dipakai sebagai encoder dalam protokol transfer. Varian ini berbeda dari RUN-original dalam empat hal: (a) causal discovery dan PageRank dihilangkan; (b) backbone di-pre-train satu kali pada dataset sumber, bukan per kasus; (c) scoring head terawasi dilatih pada label target; dan (d) DLinear, yang sebenarnya adalah model peramalan, dipakai sebagai encoder representasi. Adaptasi ini diperlukan karena transfer tidak dapat dipelajari dengan pipeline yang tidak punya parameter yang dapat ditransfer maupun tahap berbasis label. Hasilnya menggambarkan gagasan pre-training kontrastif RUN, dan tidak pernah dilaporkan sebagai kinerja RUN. \todo{Tetapkan keluaran DLinear mana yang dipakai sebagai vektor representasi. Konsekuensi: pilihan ini menentukan apakah RUN-encoder merepresentasikan deret masukan atau hasil peramalan, sehingga memengaruhi hasil C2 dan C3; keputusan perlu divalidasi sebelum eksperimen dijalankan.}
Tabel~\ref{tab:methods} merangkum mekanisme ketiga metode, dan Bagian~\ref{sec:fair} menjelaskan cara membandingkannya.
\begin{table}
\caption{Mekanisme ketiga metode yang dibandingkan, sebagaimana dijelaskan pada paper aslinya.}\label{tab:methods}
\centering
\small
\begin{tabular}{p{1.9cm}p{2.9cm}p{3.1cm}p{3.2cm}}
\toprule
 & TS2Vec \cite{yue2022ts2vec} & TimeMAE \cite{cheng2026timemae} & RUN \cite{lin2024run}\\
\midrule
Paradigma & SSL kontrastif & SSL masked autoencoder & Neural Granger causal discovery dengan pre-training kontrastif\\
Backbone & CNN dilated (10 blok residual) & Transformer encoder dengan token sub-series dan decoupled encoder terpisah & DLinear (model peramalan)\\
Sinyal pretext & Kontras temporal dan antar-instance secara hierarkis antar potongan deret & Masked codeword classification dan regresi target momentum (masking 60\%) & Kesepakatan dua potongan deret, tanpa pasangan negatif\\
Memakai label? & Hanya saat fine-tuning & Hanya saat fine-tuning & Tidak (per kasus, tak terawasi)\\
Bagian yang dapat ditransfer & Encoder & Encoder & Tidak ada pada RUN-original; backbone pada RUN-encoder\\
Keluaran untuk RCA & Representasi $\rightarrow$ scoring head & Representasi $\rightarrow$ scoring head & Graf kausal $\rightarrow$ rangking PageRank\\
\bottomrule
\end{tabular}
\end{table}


\subsection{Rangking Tingkat Service untuk Semua Metode}\label{sec:runmap}
AC@$k$ hanya dapat dibandingkan antar metode jika setiap metode merangking service kandidat yang sama. Heuristik, TS2Vec, TimeMAE, dan RUN-encoder memberi skor langsung pada service. RUN-original merangking node KPI, sehingga keluarannya dipetakan ke service. Node pada graf RUN adalah channel yang tersedia ($M_{i,j,c}=1$) dari service kandidat, dan skor sebuah service adalah skor PageRank maksimum di antara node-nya. Ini setara dengan menempatkan sebuah service pada posisi KPI pertamanya dalam rangking KPI. Akibatnya, $k$ KPI teratas pada RUN dapat mencakup kurang dari $k$ service, sehingga AC@$k$ tingkat service berbeda dari HR@$k$ tingkat KPI yang dilaporkan di \cite{lin2024run} dan tidak boleh dibandingkan langsung dengannya. \todo{Pemetaan KPI ke service (skor maksimum) adalah pilihan rancangan yang belum divalidasi. Konsekuensi: aturan lain (misalnya rata-rata skor KPI) dapat mengubah rangking RUN-original dan karenanya perbandingan C1; validasi aturan ini dengan tim atau pada data validasi sebelum eksperimen dijalankan.} RUN membutuhkan sebuah KPI pemicu (trigger), padahal kasus RCAEval tidak menyediakannya. Aturan trigger adalah keputusan metodologis yang memengaruhi RUN-original, karena trigger yang dipilih dari data pasca-injeksi memakai informasi yang sama dengan heuristik. \todo{Tetapkan aturan trigger (misalnya KPI dengan rata-rata $|z|$ terbesar setelah injeksi, atau KPI latency frontend yang tetap), pilih hanya pada kasus validasi, dan laporkan. Konsekuensi: aturan data-driven membuat RUN-original dan heuristik berbagi informasi yang sama sehingga keunggulan atau kekalahan RUN sulit ditafsirkan; aturan tetap bergantung pada pengetahuan topologi sistem.} Aturan skor seri pada Bagian~\ref{sec:metrics} diterapkan setelah pemetaan.

\subsection{Skenario Pembelajaran Transfer}
Misalkan $\mathcal{S}$ adalah dataset sumber untuk pre-training dan $\mathcal{T}$ dataset target untuk fine-tuning dan pengujian.
\begin{itemize}
\item \textbf{EX-1 (in-domain).} $\mathcal{S}=\mathcal{T}\in\{\mathrm{RE1},\mathrm{RE2},\mathrm{RE3}\}$.
\item \textbf{EX-2 (transfer).} $\mathcal{S}\neq\mathcal{T}$, dengan $\mathcal{T}\in\{\mathrm{RE1},\mathrm{RE2}\}$ dan $\mathcal{S}$ dataset RCAEval lainnya. RE3 tidak dipakai sebagai target karena pada eksperimen TS2Vec setiap model sudah mencapai AC@1 sebesar 1,00 di dataset ini, sehingga RE3 tidak dapat membedakan antar-encoder.
\end{itemize}
Pada EX-2, tiga encoder dibandingkan pada target yang sama: encoder \emph{transfer} (di-pre-train pada $\mathcal{S}$), encoder \emph{in-domain} (di-pre-train pada $\mathcal{T}$, yaitu encoder hasil pre-training pada EX-1), dan encoder \emph{acak}.

\subsection{Strategi Fine-tuning dan Kontrol}
Scoring head $h$ memetakan representasi setiap service menjadi satu skalar, sehingga $s_{i,j}=h\bigl(g(\mathbf{X}_{i,j},\mathbf{M}_{i,j})\bigr)$. \todo{Konfirmasi arsitektur head, cara mask masuk ke encoder, dan loss pelatihan (misalnya cross-entropy atas $N_i$ service). Konsekuensi: rincian ini harus identik untuk semua encoder yang dapat memakainya; jika berbeda, perbandingan C2 dan C3 tidak lagi mengisolasi efek pre-training.} Mengikuti istilah TimeMAE \cite{cheng2026timemae}:
\begin{itemize}
\item \textbf{FineLast:} $g$ dibekukan dan hanya $h$ yang dilatih, sehingga hasilnya mencerminkan kualitas representasi.
\item \textbf{FineAll:} $g$ dan $h$ dilatih bersama, sehingga seluruh parameter beradaptasi dengan target.
\end{itemize}
Data berlabel dipakai pada fraksi 10\%, 25\%, 50\%, dan 100\% dari kasus latih target. \todo{Nyatakan cara pengambilan subset dan apakah subset bersarang (nested). Konsekuensi: subset yang tidak bersarang menambah variansi antar-fraksi, sehingga tren terhadap jumlah label dapat bercampur dengan efek pemilihan sampel.} Dua kontrol dipakai. \emph{Encoder acak} memiliki arsitektur dan mode fine-tuning yang sama tetapi tanpa pre-training; pre-training dianggap bermanfaat hanya jika mengalahkan kontrol ini. \emph{Heuristik} tidak memerlukan pelatihan: skor sebuah service adalah rata-rata $|z|$ pada 240\,detik setelah injeksi dan pada channel yang tersedia,
\begin{equation}
s_{i,j}=\frac{\sum_{t\in\text{pasca}}\sum_{c=1}^{C}M_{i,j,c}\,|z_{i,j,t,c}|}{|\text{pasca}|\sum_{c=1}^{C}M_{i,j,c}},\qquad |\text{pasca}|=240 .
\end{equation}

\subsection{Metrik Evaluasi}\label{sec:metrics}
Untuk kasus $i$, rangking akar masalah adalah
\begin{equation}
r_i = 1 + \bigl|\{\,j\neq c_i : s_{i,j}\ge s_{i,c_i}\,\}\bigr|,
\end{equation}
yang menempatkan service dengan skor seri terhadap akar masalah di atasnya, sehingga skor seri tidak dapat menaikkan akurasi (aturan seri pesimistis). \todo{Aturan seri pesimistis belum divalidasi terhadap kode eksperimen TS2Vec yang sudah berjalan. Konsekuensi: jika kode lama memakai aturan lain (misalnya urutan indeks atau rangking acak), angka heuristik, RUN-original, dan AC@k lain harus dihitung ulang dengan aturan yang sama agar dapat dibandingkan.} Untuk himpunan uji $\mathcal{D}$,
\begin{equation}
\mathrm{AC@}k=\frac{1}{|\mathcal{D}|}\sum_{i\in\mathcal{D}}\mathbb{1}[r_i\le k],\qquad
\mathrm{Avg@}5=\frac{1}{5}\sum_{k=1}^{5}\mathrm{AC@}k .
\end{equation}
AC@$k$ identik dengan Top-$K$ pada ART dan HR@$k$ pada RUN jika item yang dirangking sama, dan Avg@5 adalah definisi pada ART (Bagian 5.1.3 dari \cite{sun2024art}). AC@1 adalah metrik utama dan AC@3 metrik sekunder; Avg@5 dilaporkan hanya jika diterapkan secara konsisten pada semua metode. Karena jumlah service per sistem kecil, tingkat kebetulan (chance level) untuk rangking acak seragam adalah
\begin{equation}
\mathrm{Chance@}k=\frac{1}{|\mathcal{D}|}\sum_{i\in\mathcal{D}}\frac{\min(k,N_i)}{N_i},\qquad
\mathrm{Chance\ Avg@}5=\frac{1}{5}\sum_{k=1}^{5}\mathrm{Chance@}k ,
\end{equation}
dan $\mathrm{AC@}k=1$ berlaku secara trivial ketika $N_i\le k$. Karena itu, chance level dilaporkan di samping AC@3 dan Avg@5. Selisih terhadap kontrol diuji dengan paired bootstrap atas kasus uji (10.000 resampel). Uji ini tidak menangkap variansi antar-seed; karena banyak perbandingan dilakukan, pengujian diperlakukan sebagai eksploratori kecuali dikenakan koreksi perbandingan ganda. Semua hasil adalah rata-rata dan simpangan baku dari tiga seed acak.

\section{Dataset dan Pengaturan Eksperimen}\label{sec:setup}
\subsection{Dataset}
Eksperimen memakai tiga dataset RCAEval \cite{pham2024rcaeval}, masing-masing mencakup sistem Online Boutique, Sock Shop, dan Train Ticket (Tabel~\ref{tab:datasets}). RE1 hanya berisi metrik, sedangkan RE2 dan RE3 juga berisi log dan trace; penelitian ini hanya memakai channel metrik yang dicantumkan pada Bagian~\ref{sec:method}, sehingga ketiga dataset dapat dibandingkan. RE1 dan RE2 memuat fault sumber daya dan jaringan (RE2 menambahkan fault socket), sedangkan RE3 memuat fault tingkat kode. Setelah pembagian 60/20/20, data uji berisi 75 kasus (RE1), 45 kasus (RE2), dan 18 kasus (RE3). Dataset D1 dari ART tidak dipakai pada versi penelitian ini.
\begin{table}
\caption{Dataset RCAEval yang dipakai dalam penelitian ini \cite{pham2024rcaeval}. Setiap dataset mencakup Online Boutique, Sock Shop, dan Train Ticket.}\label{tab:datasets}
\centering
\small
\begin{tabular}{lrlll}
\toprule
Dataset & Kasus & Kasus per sistem & Telemetri & Tipe fault\\
\midrule
RE1 & 375 & 125 & Hanya metrik & CPU, MEM, DISK, DELAY, LOSS\\
RE2 & 270 & 90 & Metrik, log, trace & Fault RE1 dan SOCKET\\
RE3 & 90 & 30 & Metrik, log, trace & Fault tingkat kode (F1--F5)\\
\bottomrule
\end{tabular}
\end{table}

\todo{Tambahkan jumlah service aplikasi per sistem setelah preprocessing, interval sampling asli tiap dataset, dan jumlah kasus yang kehilangan setiap channel. Konsekuensi: tanpa jumlah service, chance level AC@3 dan Avg@5 tidak dapat dihitung dan AC@3 tidak dapat ditafsirkan.}

\subsection{Rancangan Eksperimen}
Tabel~\ref{tab:config} mencantumkan konfigurasi eksperimen. EX-1 melatih dan menguji pada dataset yang sama. EX-2 melakukan pre-training pada dataset yang berbeda, sehingga setiap target dievaluasi dengan encoder transfer, in-domain, dan acak. Untuk encoder yang dapat dilatih (TS2Vec, TimeMAE, RUN-encoder), setiap konfigurasi disilangkan dengan dua mode fine-tuning dan empat fraksi label, lalu diulang dengan tiga seed pada satu pembagian data tetap.
\begin{table}
\caption{Konfigurasi eksperimen; $\mathcal{S}$ adalah dataset sumber dan $\mathcal{T}$ dataset target. Encoder in-domain pada EX-2 adalah encoder hasil pre-training pada EX-1.}\label{tab:config}
\centering
\small
\begin{tabular}{p{1.9cm}p{3.8cm}p{5.3cm}}
\toprule
 & EX-1 (in-domain) & EX-2 (transfer)\\
\midrule
Pre-training & $\mathcal{S}=\mathcal{T}$ & $\mathcal{S}\neq\mathcal{T}$, $\mathcal{S}\in\{\mathrm{RE1},\mathrm{RE2},\mathrm{RE3}\}$\\
Fine-tuning dan uji & $\mathcal{T}\in\{\mathrm{RE1},\mathrm{RE2},\mathrm{RE3}\}$ & $\mathcal{T}\in\{\mathrm{RE1},\mathrm{RE2}\}$; RE3 bukan target (jenuh)\\
Encoder & Pre-trained, acak & Transfer, in-domain, acak\\
Mode fine-tuning & FineLast, FineAll & FineLast, FineAll\\
Fraksi label & \todo{konfirmasi} & 10\%, 25\%, 50\%, 100\%\\
Seed & 3 & 3\\
Metrik & AC@1 (AC@3, Avg@5) & AC@1 (AC@3, Avg@5); paired bootstrap, 10.000 resampel\\
\bottomrule
\end{tabular}
\end{table}


\subsection{Aturan Perbandingan}\label{sec:fair}
Mekanisme pelatihan ketiga metode berbeda, sehingga satu protokol tunggal tidak dipaksakan kepada semuanya. Eksperimen menjawab tiga pertanyaan, masing-masing dengan aturan perbandingannya sendiri.
\begin{description}
\item[C1: Kinerja RCA (metode terhadap metode).] Setiap metode dijalankan pada konfigurasi terbaiknya yang dipilih pada kasus validasi, dengan kasus, seed, data uji, dan metrik yang sama (Bagian~\ref{sec:metrics}); anggaran tuning tiap metode dilaporkan. RUN-original dan heuristik disertakan sebagai acuan tanpa transfer. Perbandingan ini mengukur apa yang akan diperoleh praktisi dari tiap metode, bukan nilai pre-training.
\item[C2: Efek pre-training dan transfer (di dalam satu metode).] Untuk setiap encoder yang dapat dilatih, encoder hasil pre-training atau transfer dibandingkan dengan encoder acak berarsitektur sama pada head, label, dan mode fine-tuning yang sama, menggunakan paired bootstrap. AC@1 absolut tidak dibandingkan antar metode dalam analisis ini; hanya keuntungan setiap metode terhadap encoder acaknya sendiri yang dibandingkan, dan dinyatakan sebagai efek relatif.
\item[C3: Efek label dan fine-tuning (di dalam satu metode).] AC@1 dilaporkan untuk setiap fraksi label dan mode fine-tuning. RUN-original tidak memakai label dan hanya muncul sebagai garis acuan horizontal.
\end{description}

\subsection{Detail Implementasi}
Tabel~\ref{tab:impl} mencantumkan hiperparameter yang dilaporkan pada paper asli dan nilai yang dipakai di sini. Head, optimizer, aturan early stopping, dan seed dibuat sama pada seluruh kondisi untuk satu metode.
\begin{table}
\caption{Hiperparameter. Baris ``paper'' adalah nilai yang dilaporkan pada paper asli untuk tolok ukurnya masing-masing dan belum tentu sama dengan yang dipakai di sini; nilai ``dipakai'' harus dilengkapi untuk penelitian ini.}\label{tab:impl}
\centering
\scriptsize
\setlength{\tabcolsep}{2.5pt}
\begin{tabular}{p{1.9cm}p{3.0cm}p{3.2cm}p{2.5cm}}
\toprule
Pengaturan & TS2Vec \cite{yue2022ts2vec} & TimeMAE \cite{cheng2026timemae} & RUN \cite{lin2024run}\\
\midrule
Backbone (paper) & Proyeksi masukan, CNN dilated dengan 10 blok residual & Transformer, 8 layer, 4 head, dropout 0,2; decoupled encoder 6 layer & DLinear\\
Ukuran representasi (paper) & 320 (klasifikasi) & 64 & tidak disebutkan\\
Spesifikasi pre-training (paper) & Cropping; timestamp masking dengan $p=0{,}5$ & Rasio mask 60\%; panjang sub-series dicari pada $\{4,8,12,16\}$; ukuran codebook dicari pada $\{64,\dots,512\}$; $\alpha=1$, $\beta$ dicari pada $\{1,\dots,10\}$; momentum 0,99 & Window 32; 50 epoch untuk pre-training dan fine-tuning\\
Optimizer (paper) & Tidak ada di teks utama (appendix) & Adam, lr 0,001, batch 64 & Adam, lr 0,001, batch 128\\
Dipakai pada penelitian ini & \todo{} & \todo{} & \todo{}\\
\midrule
Head, loss, epoch, early stopping & \multicolumn{3}{l}{\todo{disamakan antar metode sejauh encoder memungkinkan}}\\
Perangkat keras, waktu komputasi & \multicolumn{3}{l}{\todo{}}\\
\bottomrule
\end{tabular}
\end{table}


\section{Hasil Eksperimen}\label{sec:results}
Entri bertanda \pending\ adalah eksperimen yang belum dijalankan; tidak ada nilai yang diperkirakan. Hasil yang tersedia saat ini hanya untuk TS2Vec dan berasal dari laporan kemajuan proyek. Metrik yang dipakai adalah AC@1 kecuali dinyatakan lain.

\subsection{C1: Kinerja RCA Antar Metode}
Tabel~\ref{tab:compare} adalah perbandingan C1 yang didefinisikan pada Bagian~\ref{sec:fair}.
\begin{table}
\caption{Perbandingan C1: kinerja RCA setiap metode pada konfigurasi terbaiknya yang dipilih pada kasus validasi (AC@1 dan AC@3 pada data uji, rata-rata dari 3 seed). RUN-encoder adalah varian adaptasi, bukan RUN asli. Semua sel masih terbuka sampai eksperimen dijalankan; nilai TS2Vec pada Tabel~\ref{tab:ex1} berasal dari konfigurasi tetap dan belum mengikuti aturan seleksi ini.}\label{tab:compare}
\centering
\small
\begin{tabular}{lcccc}
\toprule
 & \multicolumn{2}{c}{RE1} & \multicolumn{2}{c}{RE2}\\
\cmidrule(lr){2-3}\cmidrule(lr){4-5}
Metode & AC@1 & AC@3 & AC@1 & AC@3\\
\midrule
Chance level & \pending & \pending & \pending & \pending\\
Heuristik & 0,81 & \pending & 0,80 & \pending\\
RUN-original (tanpa transfer) & \pending & \pending & \pending & \pending\\
TS2Vec & \pending & \pending & \pending & \pending\\
TimeMAE & \pending & \pending & \pending & \pending\\
RUN-encoder & \pending & \pending & \pending & \pending\\
\bottomrule
\end{tabular}
\end{table}

\todo{Isi setelah eksperimen TimeMAE dan RUN selesai. Laporkan rata-rata $\pm$ SD dari 3 seed. Kerangka interpretasi: bandingkan metode hanya di dalam tabel ini; nyatakan anggaran tuning; jangan merangking metode dari hasil parsial.}

\subsection{EX-1: Pre-training dan Evaluasi In-domain}
Tabel~\ref{tab:ex1} melaporkan pre-training dan pengujian pada dataset yang sama.
\begin{table}
\caption{EX-1: AC@1 pada data uji, rata-rata $\pm$ SD dari 3 seed (pre-training dan pengujian pada dataset yang sama). Nilai TS2Vec berasal dari slide proyek; \pending\ menandai eksperimen yang belum dijalankan.}\label{tab:ex1}
\centering
\scriptsize
\setlength{\tabcolsep}{3pt}
\begin{tabular}{llccccc}
\toprule
 & & & \multicolumn{2}{c}{FineLast} & \multicolumn{2}{c}{FineAll}\\
\cmidrule(lr){4-5}\cmidrule(lr){6-7}
Dataset (kasus uji) & Metode & Heuristik & Pre-trained & Acak & Pre-trained & Acak\\
\midrule
RE1 (75) & TS2Vec  & 0,81 & 0,92$\pm$0,02 & 0,94$\pm$0,01 & 0,95$\pm$0,01 & 0,93$\pm$0,02\\
         & TimeMAE & 0,81 & \pending & \pending & \pending & \pending\\
         & RUN-encoder & 0,81 & \pending & \pending & \pending & \pending\\
\midrule
RE2 (45) & TS2Vec  & 0,80 & 0,89$\pm$0,02 & 0,86$\pm$0,01 & 0,86$\pm$0,01 & 0,87$\pm$0,02\\
         & TimeMAE & 0,80 & \pending & \pending & \pending & \pending\\
         & RUN-encoder & 0,80 & \pending & \pending & \pending & \pending\\
\midrule
RE3 (18) & TS2Vec  & 0,67 & 1,00$\pm$0,00 & 1,00$\pm$0,00 & 1,00$\pm$0,00 & 1,00$\pm$0,00\\
         & TimeMAE & 0,67 & \pending & \pending & \pending & \pending\\
         & RUN-encoder & 0,67 & \pending & \pending & \pending & \pending\\
\bottomrule
\end{tabular}
\end{table}

Pada TS2Vec, model hasil fine-tuning melampaui heuristik pada RE1 dan RE2 di keempat konfigurasi, dan encoder acak sama baiknya dengan encoder hasil pre-training pada sebagian besar sel; selisih terbesar antara pre-training dan acak adalah 0,03, setara satu atau dua kasus uji. Pada RE3, setiap konfigurasi mencapai 1,00. \todo{Tambahkan baris TimeMAE dan RUN-encoder setelah eksperimen dijalankan.}

\paragraph{Kontrol penghapusan anomali.} Data 240\,detik setelah injeksi diganti dengan salinan 240\,detik sebelum injeksi (Tabel~\ref{tab:anomaly}). AC@1 turun menjadi 0,03, 0,09, dan 0 pada RE1, RE2, dan RE3, jauh di bawah nilai aslinya. Hal ini menunjukkan bahwa model bergantung pada anomali dan tidak hanya pada identitas service. Slide melaporkan acuan tebakan acak sebesar $\pm$0,06; arti pastinya harus diklarifikasi sebelum dipakai sebagai chance level (Bagian~\ref{sec:metrics}).
\begin{table}
\caption{Kontrol penghapusan anomali (TS2Vec): AC@1 dengan data asli dan dengan jendela pasca-injeksi yang diganti salinan data pra-injeksi. Nilai berasal dari slide proyek, yang tidak menyebutkan konfigurasi model di baliknya.}\label{tab:anomaly}
\centering
\small
\begin{tabular}{lccc}
\toprule
Dataset & Data asli & Anomali dihapus & Tebakan acak (slide)\\
\midrule
RE1 & 0,91 & 0,03 & $\pm$0,06\\
RE2 & 0,9  & 0,09 & $\pm$0,06\\
RE3 & 1    & 0    & $\pm$0,06\\
\bottomrule
\end{tabular}
\end{table}


\subsection{EX-2: Transfer Lintas Dataset}
Tabel~\ref{tab:ex2} dan Gambar~\ref{fig:ex2} melaporkan RE2 sebagai target dengan mode FineLast, sedangkan Tabel~\ref{tab:bootstrap} melaporkan paired bootstrap terhadap encoder acak.
\begin{table}
\caption{EX-2: AC@1 pada RE2 (45 kasus uji) dengan FineLast, menurut persentase data latih berlabel, rata-rata dari 3 seed. Nilai TS2Vec dibaca dari grafik slide (dua desimal; SD tidak dilaporkan di sana). Pengaturan 10\% tidak memiliki nilai pada slide (``semua encoder setara''); tl = tidak dilaporkan. Heuristik memperoleh 0,80.}\label{tab:ex2}
\centering
\small
\begin{tabular}{llcccc}
\toprule
Metode & Encoder & 10\% & 25\% & 50\% & 100\%\\
\midrule
TS2Vec & Transfer dari RE1 & tl & 0,90 & 0,91 & 0,91\\
       & Transfer dari RE3 & tl & 0,86 & 0,88 & 0,87\\
       & In-domain (RE2)   & tl & 0,87 & 0,90 & 0,89\\
       & Acak              & tl & 0,84 & 0,84 & 0,86\\
\midrule
TimeMAE & Transfer dari RE1 / RE3, in-domain, acak & \pending & \pending & \pending & \pending\\
RUN-encoder & Transfer dari RE1 / RE3, in-domain, acak & \pending & \pending & \pending & \pending\\
\bottomrule
\end{tabular}
\end{table}

\begin{figure}
\centering
\begin{tikzpicture}
\begin{axis}[
  ybar, bar width=7pt, width=\textwidth, height=5.2cm,
  ymin=0.7, ymax=1.0, ylabel={AC@1 (data uji RE2)},
  yticklabel style={/pgf/number format/.cd, fixed, precision=1, use comma},
  scaled y ticks=false,
  symbolic x coords={p25,p50,p100}, xtick=data,
  xticklabels={25\%,50\%,100\%},
  xlabel={Data latih berlabel},
  legend style={at={(0.5,1.05)},anchor=south,legend columns=-1,font=\scriptsize},
  tick label style={font=\scriptsize}, label style={font=\scriptsize},
  enlarge x limits=0.25]
\addplot coordinates {(p25,0.90) (p50,0.91) (p100,0.91)};
\addplot coordinates {(p25,0.86) (p50,0.88) (p100,0.87)};
\addplot coordinates {(p25,0.87) (p50,0.90) (p100,0.89)};
\addplot coordinates {(p25,0.84) (p50,0.84) (p100,0.86)};
\legend{Dari RE1, Dari RE3, In-domain (RE2), Acak}
\end{axis}
\end{tikzpicture}
\caption{EX-2, TS2Vec, FineLast: AC@1 pada RE2 menurut fraksi data latih berlabel (nilai dari slide proyek, rata-rata 3 seed). Heuristik memperoleh 0,80 pada RE2; sumbu-y dimulai dari 0,7. Seri TimeMAE dan RUN-encoder akan ditambahkan ketika eksperimennya tersedia.}\label{fig:ex2}
\end{figure}

\begin{table}
\caption{Paired bootstrap (10.000 resampel) untuk selisih AC@1 antara encoder transfer dan encoder acak, hanya TS2Vec. Signifikansi sesuai pernyataan pada slide; interval kepercayaan hanya ditampilkan secara grafis di slide dan tidak direproduksi di sini.}\label{tab:bootstrap}
\centering
\small
\begin{tabular}{llcc}
\toprule
Target (sumber) & Pengaturan & Selisih & Signifikan\\
\midrule
RE2 (RE1) & FineLast 10\%  & 0,00 & tidak\\
          & FineLast 25\%  & +0,06 & ya\\
          & FineLast 50\%  & +0,07 & ya\\
          & FineLast 100\% & +0,05 & ya\\
          & FineAll 25\%   & +0,03 & tidak\\
          & FineAll 50\%   & +0,04 & ya\\
\midrule
RE1 (RE2) & FineLast 10\%  & 0,00 & tidak\\
          & FineLast 25\%  & $-$0,02 & tidak\\
          & FineLast 50\%  & $-$0,01 & tidak\\
          & FineLast 100\% & 0,00 & tidak\\
\bottomrule
\end{tabular}
\end{table}

Pada TS2Vec dengan RE2 sebagai target, encoder dari RE1 menghasilkan AC@1 sekitar 0,90--0,91 dengan 25\%--100\% label, dibandingkan 0,84--0,86 untuk encoder acak, dan selisihnya signifikan pada 4 dari 6 pengaturan yang dijalankan. Dengan 10\% label, semua encoder setara. Encoder dari RE3 tidak memberi manfaat. Dengan RE1 sebagai target, tidak ada encoder sumber yang memberi selisih signifikan.

\subsection{C2 dan C3: Pre-training, Label, dan Fine-tuning}
Tabel~\ref{tab:ex1}--\ref{tab:bootstrap} adalah bukti untuk kedua pertanyaan ini. Pada TS2Vec, FineAll tidak lebih baik daripada FineLast pada pengaturan yang dilaporkan. Untuk TimeMAE dan RUN-encoder, tabel yang sama dan Gambar~\ref{fig:ex2} akan diperluas, dan keuntungan terhadap encoder acak milik masing-masing metode adalah besaran yang dibandingkan antar metode.

\subsection{Metrik Tambahan}
\todo{AC@3 dan Avg@5 beserta chance level, setelah jumlah service per sistem diketahui dan kode metrik diterapkan pada semua metode.}

\section{Pembahasan}\label{sec:discussion}
\subsection{Memisahkan Manfaat Fine-tuning dari Manfaat Pre-training}
Pada TS2Vec, keunggulan atas heuristik ada baik pada encoder acak maupun pada encoder hasil pre-training di EX-1, sehingga keunggulan tersebut tidak dapat dikaitkan dengan pre-training. Salah satu penjelasan yang masuk akal adalah bahwa encoder konvolusional yang diinisialisasi acak sudah menyediakan fitur yang berguna ketika sebuah head dilatih di atasnya, tetapi penjelasan ini merupakan hipotesis dan belum diuji di sini. \todo{Putuskan apakah akan menambah eksperimen (misalnya linear probe pada fitur mentah) untuk menguji hipotesis ini. Konsekuensi: tanpa uji tersebut, penjelasan hanya boleh ditulis sebagai dugaan.}

\subsection{Kapan Transfer Membantu dan Kapan Tidak}
Satu-satunya efek transfer yang teramati sejauh ini adalah RE1 $\rightarrow$ RE2: efeknya kecil dan tidak terlihat pada 10\% label, dengan RE3 sebagai sumber, maupun dengan RE1 sebagai target. RE1 dan RE2 memiliki tipe fault yang sama kecuali fault socket, sehingga representasi RE1 mungkin sesuai untuk RE2, sedangkan RE3 memuat fault tingkat kode. Penjelasan berbasis kesamaan tipe fault ini adalah hipotesis yang perlu diperiksa, misalnya dengan membandingkan hasil per tipe fault. \todo{Laporkan hasil per tipe fault jika data memungkinkan. Konsekuensi: tanpa prediksi per kasus yang tersimpan, analisis ini membutuhkan eksperimen ulang.} Mode FineAll tidak lebih baik daripada FineLast; laporan kemajuan mengaitkannya dengan overfitting yang cepat, dan hal ini perlu diverifikasi dengan kurva pelatihan.

\subsection{Membaca Perbandingan Antar Metode}
\todo{Setelah eksperimen: (i) laporkan C1 lebih dulu sebagai perbandingan praktis; (ii) laporkan keuntungan setiap metode terhadap encoder acaknya sendiri (C2) dan nyatakan apakah perbedaannya melampaui ketidakpastiannya; (iii) bahas apakah perbedaan berasal dari tugas pretext, backbone, atau anggaran pre-training; (iv) bahas RUN-encoder secara terpisah dari RUN-original dan jangan menyimpulkan apa pun tentang RUN dari RUN-encoder.}
Aturan keputusan yang ditetapkan sebelumnya: pre-training dinyatakan bermanfaat hanya jika paired bootstrap terhadap encoder acak pada mode yang sama signifikan, dan dinyatakan tidak memiliki efek yang terdeteksi jika tidak.

\subsection{Keterbatasan}
Data uji kecil (75, 45, dan 18 kasus), sehingga selisih satu atau dua kasus berada dalam rentang derau. Pembagian data hanya satu dan seed berjumlah tiga. Bootstrap hanya me-resample kasus. Waktu injeksi diasumsikan diketahui. Heuristik dan semua encoder hanya memakai lima channel metrik. RE3 jenuh sehingga tidak informatif. Hasil saat ini hanya untuk TS2Vec, sehingga belum ada kesimpulan tentang kelebihan relatif ketiga metode. Generalisasi lintas dataset masih terbuka karena transfer hanya membantu pada satu arah.

\section{Kesimpulan dan Penelitian Lanjutan}\label{sec:conclusion}
\todo{Tulis setelah seluruh eksperimen selesai. Struktur: (1) apa yang dikaji; (2) jawaban RQ1--RQ3 masing-masing satu kalimat disertai angka; (3) keterbatasan; (4) penelitian lanjutan. Jangan menyatakan rangking metode sebelum hasil TimeMAE dan RUN-encoder tersedia.}

Penelitian lanjutan yang sudah direncanakan meliputi penyelarasan fitur antar dataset dengan PCA, evaluasi pada dataset D1 dari ART, dan penambahan pendekatan ART sebagai metode acuan.

\begin{credits}
\subsubsection{Pernyataan Kepentingan.}
\todo{Penulis perlu menyediakan pernyataan konflik kepentingan.}
\end{credits}


\bibliographystyle{splncs04}
\bibliography{references}

\end{document}
